\documentclass[11pt]{article}
\usepackage[margin=1in]{geometry}
\usepackage{lmodern}
\usepackage{hyperref}
\usepackage{enumitem}
\usepackage{amsmath,amssymb,amsthm,mathtools}
\usepackage[T1]{fontenc}
\usepackage[utf8]{inputenc}
\title{\textbf{Theory-mode report: Strong-disorder expansion of the root-averaged density of states for the Anderson model on the Bethe lattice}}
\author{Generated from Scholar Inbox digest}
\date{Digest date: 2026-05-04}
\begin{document}
\maketitle
\section*{Paper information}
\begin{itemize}[leftmargin=1.5em]
\item \textbf{Title:} Strong-disorder expansion of the root-averaged density of states for the Anderson model on the Bethe lattice
\item \textbf{Author:} Masahiro Kaminaga
\item \textbf{arXiv:} \href{https://arxiv.org/abs/2605.00478}{2605.00478}
\end{itemize}
\section*{Executive summary}
This paper studies the Anderson model on the Bethe lattice in the strong-disorder regime and focuses on the \emph{root-averaged density of states}, i.e. the spectral measure seen from the root rather than a finite-volume eigenvalue counting measure. The main theorem is a local regularity and asymptotic expansion result on scaled energy windows of size $\lambda$, where $\lambda$ is the disorder strength. Under compact support and local analyticity assumptions on the single-site distribution, the averaged diagonal resolvent admits a holomorphic continuation to a complex neighborhood of a target interval once $\lambda$ is large enough. By Stieltjes inversion, the corresponding density of states on the window $E=\lambda\xi$ is absolutely continuous, real analytic in $\xi$, and has a finite-order asymptotic expansion in inverse powers of $\lambda$.

The structural point is that the leading coefficient is just the local density of the single-site potential law, while all odd coefficients vanish. The even coefficients are computable from short closed walks on the tree. So the paper turns the strong-disorder regime into a clean combinatorial-plus-complex-analytic expansion problem.
\section*{Model and observable}
The operator is the Anderson Hamiltonian on the Bethe lattice,
\[
H_\lambda = T + \lambda V,
\]
where $T$ is the adjacency operator on the regular tree and $V$ is an i.i.d. random potential. The object of interest is the root spectral measure, or equivalently the averaged diagonal Green function
\[
G_\lambda(z)=\mathbb E\,\langle \delta_o,(H_\lambda-z)^{-1}\delta_o\rangle.
\]
The density of states here means the measure recovered from $G_\lambda$ by the Stieltjes inversion formula. This is different from the more common finite-volume density of states and is better adapted to the rooted tree geometry.

The paper studies the scaled energy variable $E=\lambda \xi$ with $\xi$ in a fixed compact interval $I$ contained in a slightly larger interval $I^{\sharp}$ on which the single-site density is assumed analytic. Intuitively, when disorder is large, the on-site randomness should dominate the hopping operator, so the density of states should look like the potential density plus controlled corrections coming from short excursions away from the root.
\section*{Main results}
The main results are Theorems 2.1 and 2.2.
\begin{itemize}[leftmargin=1.5em]
\item Theorem 2.1 gives holomorphic continuation of a scaled averaged diagonal resolvent to a complex neighborhood of $I$ for all sufficiently large $\lambda$.
\item Theorem 2.2 converts this into a strong-disorder expansion for the root-averaged density of states on the window $\lambda I$.
\end{itemize}
More concretely, if $\rho_\lambda(\lambda\xi)$ denotes the density in scaled form, then the paper proves an expansion of the form
\[
\rho_\lambda(\lambda\xi)
= r(\xi)+\sum_{m=1}^{M} \lambda^{-2m} a_{2m}(\xi)+O(\lambda^{-2M-2}),
\]
uniformly on the interval of interest. Here $r(\xi)$ is the density of the single-site law and the odd powers of $\lambda^{-1}$ are absent. This vanishing of odd coefficients is not cosmetic; it comes from the parity structure of closed walks on the tree and is one of the cleanest outputs of the method.

For the uniform single-site distribution, the first nonzero correction term is computed explicitly. So the result is not just abstract asymptotics: it produces concrete formulas.
\section*{Why the Bethe lattice helps}
The tree geometry removes the combinatorial mess of arbitrary loops in finite-dimensional lattices. Every return to the root can be organized by a closed walk profile, and the adjacency operator contributes through walk lengths and multiplicities. In the strong-disorder regime, the hopping part acts perturbatively, and the Green function can be expanded into a random-walk series around the root.

The crucial advantage is that on a tree, short closed walks have a rigid structure. That makes it possible to express coefficients of the resolvent expansion as finite sums over occupation profiles of such walks. The coefficients are therefore local, combinatorial, and explicit in principle.
\section*{Proof strategy}
The argument has three layers.
\subsection*{1. Random-walk expansion}
The resolvent at the root is expanded using the walk representation of powers of the adjacency operator. After scaling by $\lambda$, one obtains a formal series in $\lambda^{-1}$ whose coefficients depend on products of single-site Stieltjes transforms evaluated along vertices visited by a closed walk. The tree structure allows one to organize these contributions by occupation multiplicities.

A key combinatorial fact, reflected in Lemma 3.1, is that odd orders vanish. This fits the intuition that a closed walk on a bipartite tree has even length, so only even-order perturbative corrections survive in the density expansion.
\subsection*{2. Complex-analytic control of the single-site transforms}
Formal expansions are not enough; one needs analytic continuation in a full complex neighborhood of the target interval. This is where the local analyticity of the single-site density enters. The paper studies single-site Stieltjes transforms and proves their continuation and uniform control near the real axis. Lemma 4.1 supplies the analytic continuation estimates needed to sum the walk contributions safely.
\subsection*{3. Stieltjes inversion and extraction of the density}
Once the averaged diagonal resolvent is controlled holomorphically, one applies the Stieltjes inversion formula to pass from resolvent asymptotics to density asymptotics. Real analyticity of the resulting density follows from holomorphic continuation, and the finite-order expansion is then read off term by term.
\section*{What is actually new}
The novelty is not ``there is a perturbative expansion'' in some vague sense. It is the combination of three precise features:
\begin{itemize}[leftmargin=1.5em]
\item a local analyticity theorem for the scaled resolvent on real-energy windows,
\item a finite-order strong-disorder expansion for the root density of states,
\item an explicit combinatorial description of the coefficients through short closed walks.
\end{itemize}
The leading coefficient being the local potential density is exactly what one expects physically, but the theorem quantifies the correction mechanism and shows that the expansion is rigid enough to kill all odd orders.
\section*{Interpretation}
In scaled variables, the strong-disorder picture is simple: the operator is dominated by the diagonal potential, and hopping only produces rare return corrections. The density of states therefore begins by reproducing the single-site density. The next corrections are suppressed by even powers of $\lambda^{-1}$ and encode the probability-weighted effect of short excursions on the tree.

This is a nice example where disorder asymptotics are not merely probabilistic but genuinely analytic: the regularity of the potential law gets transferred to regularity of the spectral density in the strong-disorder window.
\section*{Limitations}
The assumptions are local and fairly strong. The single-site law is required to have compact support and a density that is analytic on a neighborhood of the interval of interest. So the result is not aimed at rough distributions or singular laws. Also, the observable is the root-averaged density of states on the Bethe lattice; it does not directly settle analogous expansions for finite-dimensional lattices or for the usual infinite-volume integrated density of states.
\section*{Bottom line}
A clean, well-targeted strong-disorder paper. The key message is that for the Anderson model on the Bethe lattice, the root density of states near energies $E=\lambda\xi$ is not only perturbatively close to the single-site potential density, but admits a controlled analytic expansion with explicitly structured coefficients. The mix of walk combinatorics and complex analysis is the heart of the paper.
\end{document}
